% Conclusion: five moves (claim, key idea, what was done, one number, outlook); ~115 words; no cite, no ref.
\section{Conclusion}
\label{sec:conclusion}

% sources-of-gain table (D-35; generated by scripts/make_tables.py, simulated rows, evaluation/dac_2x2 Sec. 2i):
% defined here, after the Conclusion heading, so that it lands at the top of p. 11 R and not in Fig. 9's column
% (one float per column); placement checked on the float map after each build
\begin{table}[t]
  \caption{Useful-MAC utilization from stock AiM to CoDe, one source at a time.}
  \label{tab:umac_src}
  \centering
  \footnotesize
  \setlength{\tabcolsep}{4pt}
  % generated by scripts/make_tables.py from evaluation/dac_2x2/figures/tab_umac_src_steps.csv (usrc_steps ids, simulated rows, layer, M = 8); do not edit
\begin{tabular}{@{}lccccc@{}}
\toprule
Design & \UMAC & \fMAC & $U$ & \multicolumn{2}{c@{}}{Speedup} \\
\cmidrule(l){5-6}
 & & & & step & cumul. \\
\midrule
\multicolumn{6}{@{}l}{\emph{Mobile:} Llama-2-7B, $V = 8$} \\
Stock AiM & 0.25 & 0.776 & 0.194 & -- & 1.00 \\
Row reuse only & 0.25 & 0.965 & 0.241 & 1.24 & 1.24 \\
\textbf{CoDe} & 1.00 & 0.843 & 0.843 & 3.49 & 4.35 \\
\midrule
\multicolumn{6}{@{}l}{\emph{Server:} Llama-3.1-8B, $V = 4$} \\
Stock AiM & 0.25 & 0.771 & 0.193 & -- & 1.00 \\
Row reuse only & 0.25 & 0.887 & 0.222 & 1.15 & 1.15 \\
\textbf{CoDe} & 1.00 & 0.739 & 0.739 & 3.33 & 3.83 \\
\bottomrule
\end{tabular}

\end{table}

% P:conc-summary
Skinny GEMMs keep LLM decoding memory-bound on mobile platforms and edge servers, and prior DRAM PIM paces every
multiply--accumulate (MAC) by a DRAM command. This paper argues that DRAM timing bounds column accesses, not MACs on a
latched column. The key to CoDe is to decouple MAC timing from DRAM command timing: one column command feeds many processing-unit beats over a tile
of tokens. We build CoDe on each rival's own memory. On LPDDR5-6400, CoDe matches LP-Spec-style's decoding speed
with 1.87$\times$ its throughput per logic mm$^2$; on HBM2, CoDe decodes batches 1.68$\times$ faster than P3-LLM-style. Since a
column's beats follow from its period and the clock alone, the decoupling carries over to any DRAM whose column
outlasts one beat.
